\chapter{Summary, Conclusion, and Implications}

\section{Introduction}
This concluding chapter synthesizes the theoretical arguments, empirical methodologies, econometric results, and interpretive discussions developed throughout this dissertation. The primary objective of this investigation has been to examine the value relevance of accounting information, specifically earnings per share (EPS) and book value per share (BVPS), among Nepalese commercial banks, and to evaluate whether the fundamental valuation relationships underwent a structural regime shift following the landmark commercial banking consolidation wave executed in early 2023 (mid-FY 2079/80).

The \textit{Master's Degree Dissertation Guidelines} (2019, p.~15) of the Faculty of Management, Tribhuvan University organize the concluding chapter into three sections: Summary, Conclusion, and Implications, with recommendations and areas for future research reported under Implications. This chapter follows that structure and, for clarity of presentation, separates the recommendations and future-research material into its own section. It is therefore organized into four core functional sections:
\begin{enumerate}
    \item \textbf{Summary of the Study} (Section~5.2), providing an integrated overview of the research problem, operational objectives, theoretical foundation, research design, and major empirical results.
    \item \textbf{Conclusions} (Section~5.3), articulating the insights, formal hypothesis testing decisions, and academic knowledge derived from the empirical findings regarding accounting value relevance, structural discontinuities, and capital market pricing behavior in Nepal.
    \item \textbf{Implications of the Study} (Section~5.4), delineating targeted implications partitioned across three institutional domains: regulatory authorities (NRB, SEBON, and NEPSE), commercial bank executives and boards, and equity investors and market participants.
    \item \textbf{Recommendations and Directions for Future Research} (Section~5.5), presenting practical policy recommendations alongside prospective avenues for subsequent scholarly inquiry.
\end{enumerate}

\section{Summary of the Study}

\subsection{Research Overview, Problem, and Objectives}
The valuation of common equity lies at the intersection of accounting disclosure, corporate governance, and capital market efficiency. Under the residual income valuation (RIV) framework established by \citet{ohlson1995} and \citet{feltham1995}, market equity value is formalized as a linear function of contemporaneous balance-sheet equity (book value per share, capturing cumulative resources) and operational profitability (earnings per share, capturing abnormal forward-looking earning capacity), anchored in the clean-surplus accounting relation.

In emerging and frontier capital markets such as Nepal, the empirical stability of these accounting valuation relationships is routinely challenged by institutional friction, regulatory shifts, information asymmetry, and structural corporate reorganizations. The commercial banking sector in Nepal represents the dominant institutional pillar of the Nepal Stock Exchange (NEPSE), accounting for more than 30\% of total market capitalization \citep{gurungdahal2025}. Over the past decade, Nepal Rastra Bank (NRB) pursued aggressive regulatory policies to consolidate the banking industry, culminating in the 2022--2023 consolidation wave during which major commercial banks executed large-scale horizontal amalgamations and statutory acquisitions.

Despite the landmark scale of this corporate restructuring, existing empirical research on Nepalese banking equity valuation exhibited substantial theoretical, methodological, and econometric deficiencies:
\begin{itemize}
    \item Almost all published Nepalese studies employed static, full-sample pooled or one-way random-effects regressions, implicitly assuming parameter stability and failing to test whether corporate combinations altered valuation slopes across the pre- and post-consolidation regimes \citep{chaudhari2025,chhetri2023,timalsina2026}. \citet{gurungdahal2025} estimated a bank fixed-effects model rather than a pooled regression, but likewise did not test a regime shift.
    \item Empirical models frequently introduced the circular Price-to-Earnings (P/E) ratio as an independent regressor predicting market share price (MPS), embedding the dependent variable within its own explanatory matrix and generating severe mechanical endogeneity and spurious correlation \citep{bhattarai2014,giri2024,joshi2023,pant2024}.
    \item Bhattarai (2014) measured share prices at fiscal year-end (Ashadh-end / mid-July), long before audited annual financial statements were completed and published, inducing severe look-ahead bias and violating temporal causality \citep{bhattarai2014}.
    \item Prior panel estimations routinely ignored the finite-sample constraints of small cross-sectional clusters, risking false discoveries and overstated statistical precision through inappropriate covariance estimators. This study's own cluster count is $G = 5$.
\end{itemize}

To resolve these empirical gaps, this dissertation addressed three core research questions:
\begin{itemize}
    \item Are fundamental accounting measures (EPS and BVPS) statistically associated with secondary market share prices observed on or after statutory annual reporting deadlines ($\mathit{MPS}^{RD}$)?
    \item Did the valuation associations of EPS and BVPS differ between the pre-consolidation period (FY 2076/77--2078/79) and the post-consolidation period (FY 2079/80--2081/82)?
    \item How sensitive are the estimated valuation coefficients to transition-year effects, entity-timing mismatches, operating profitability substitutions, inclusion of eligible dividend disclosures, bonus share adjustments, restated earnings, and sensitivity analysis excluding the NMB FY 2076/77 bank-year with the Kanchan acquisition timing mismatch?
\end{itemize}

Correspondingly, the study set out to: (i) evaluate the baseline value relevance of EPS and BVPS under a standardized post-reporting-deadline pricing framework, (ii) quantify structural shifts in valuation slopes across the 2023 consolidation wave using an interactive panel specification, and (iii) execute a nine-specification estimation and robustness suite (R1--R2 baselines plus seven robustness variants R3--R9) to establish the empirical boundary conditions of accounting valuation slopes in Nepal.

\subsection{Methodology and Research Design Synthesis}
To execute this investigation with sound econometric discipline, the study employed a longitudinal panel design spanning six fiscal years (FY 2076/77 through FY 2081/82, corresponding to 2019/20 through 2024/25). The target population encompassed all commercial banks licensed by Nepal Rastra Bank as Class `A' financial institutions. A purposive sample of five commercial banking series was selected:
\begin{itemize}
    \item Four amalgamated banking series representing the major corporate combinations of the early-2023 wave: Global IME Bank Limited (GBIME, absorbed Bank of Kathmandu), Himalayan Bank Limited (HBL, acquired Civil Bank), Kumari Bank Limited (KBL, absorbed NCC Bank), and Nepal Investment Mega Bank Limited (NIMB, amalgamation of Nepal Investment Bank and Mega Bank Nepal).
    \item One benchmark banking series: NMB Bank Limited (NMB), which maintained continuous corporate operations throughout the 2023 wave without participating in horizontal combinations (note: NMB acquired Kanchan Development Bank Limited in August 2020, adding approximately Nepalese rupees (NPR) 0.50 billion in new shares issued to Kanchan shareholders at the 85\% swap ratio, alongside a 13\% bonus share distribution, in FY 2077/78, while its corporate identity remained continuous through the 2023 consolidation).
\end{itemize}

To reduce temporal mismatch between accounting periods and price observations, market prices were observed on or after the statutory annual financial reporting deadline. Section 59(2) of the Banks and Financial Institutions Act, 2017 (BAFIA) legally mandates that commercial banks publish their audited annual balance sheets and profit and loss accounts within nine months into the subsequent fiscal year, operationalized at the end of Chaitra (mid-April) of the subsequent fiscal year \citep{bafia2017}. The dependent variable, $\mathit{MPS}^{RD}$, was defined as the closing market price on the first trading session on or after this statutory threshold, subject to a strict 5-day freshness rule.

Under this 5-day rule, 28 clean bank-year observations were realized out of an intended balanced panel of 30 bank-years (13 pre-consolidation and 15 post-consolidation observations). Two observations were excluded and dropped rather than imputed: HBL FY 2077/78 (trading resumed 54 days late on 6 June 2022) and NIMB FY 2077/78 (trading resumed 47 days late on 30 May 2022), both resulting from prolonged merger-related secondary market trading suspensions around the April 2022 deadline.

Econometric estimation was structured around two core models:
\begin{itemize}
    \item \textbf{Model 1 (Baseline Value Relevance):} Specified as a two-way fixed effects panel regression:
    \begin{equation}
    \mathit{MPS}^{RD}_{it} = \alpha + \mu_i + \lambda_t + \beta_E \mathit{EPS}_{it} + \beta_B \mathit{BVPS}_{it} + \varepsilon_{it}
    \end{equation}
    where $\mu_i$ captures unobserved bank fixed effects (with GBIME as reference) and $\lambda_t$ absorbs common fiscal-year macroeconomic shocks and NEPSE-wide price cycles (with FY 2076/77 as reference). Residual degrees of freedom were exact: $28 - 12 = 16$.
    \item \textbf{Model 2 (Regime Shift Interaction):} Specified to estimate structural changes across the early-2023 consolidation wave:
    \begin{equation}
    \begin{aligned}
    \mathit{MPS}^{RD}_{it} ={}& \alpha + \mu_i + \lambda_t + \beta_E \mathit{EPS}_{it} + \beta_B \mathit{BVPS}_{it} \\
    & + \delta_E (\mathit{EPS}_{it} \times \mathit{Post}_t) + \delta_B (\mathit{BVPS}_{it} \times \mathit{Post}_t) + \varepsilon_{it}
    \end{aligned}
    \end{equation}
    where $\mathit{Post}_t = 1$ for post-consolidation years (FY 2079/80--2081/82) and 0 for pre-consolidation years (FY 2076/77--2078/79). The main effect of $\mathit{Post}_t$ was fully absorbed by the fiscal-year fixed effects ($\lambda_t$). Residual degrees of freedom were exact: $28 - 14 = 14$.
\end{itemize}

Because the five-bank panel is too small to support a reliable cluster-robust correction, all inferential statistical tests ($t$-tests, confidence intervals, and the principal joint Wald $F$-test of $H_0: \delta_E = \delta_B = 0$) were conducted using the heteroskedasticity-consistent (HC3) covariance estimator of \citet{mackinnon1985} evaluated on Student's $t$ residual degrees of freedom. HC3 addresses heteroskedasticity, not within-bank correlation, so the within-bank independence assumption remains untested.

\subsection{Summary of Major Empirical Findings}
The empirical analysis yielded several substantive findings regarding the statistical properties, baseline relevance, regime stability, and sensitivity of accounting fundamentals. The underlying dataset is printed in Appendix~\ref{app:analysis_dataset}, and the full estimation outputs behind these findings are printed in Appendix~\ref{app:estimation_outputs}.
\begin{itemize}
    \item \textbf{Descriptive Shifts Across Regimes:}
    \begin{itemize}
        \item \textit{Market Price Compression:} Full sample mean $\mathit{MPS}^{RD}$ stood at NPR 254.34. Across regimes, mean share price dropped by 33.81\% from NPR 310.59 pre-consolidation to NPR 205.60 post-consolidation, while cross-sectional standard deviation contracted by 72.00\% (from 119.46 to 33.45), documenting a substantial market-wide price compression.
        \item \textit{Earnings Deterioration and Dispersion:} Mean EPS contracted by 40.13\% from NPR 17.90 pre-consolidation to NPR 10.72 post-consolidation. Its SD increased by 63.95\% (3.99 to 6.53), and its coefficient of variation rose from 22.3\% to 61.0\%.
        \item \textit{Dividend Reduction:} Mean DPS fell by 76.90\% (NPR 6.06 to NPR 1.40), and its post-consolidation median was zero, with 12 of 15 observations recording no cash dividend. The DPS SD declined 16.62\%, and its coefficient of variation rose from 58.8\% to 212.2\%.
        \item \textit{Book Value Stability:} Mean BVPS increased by 2.45\% from NPR 158.11 to NPR 161.97, while its SD declined by 15.30\% (NPR 19.10 to NPR 16.18). These descriptive changes do not identify their cause.
    \end{itemize}
    \item \textbf{Correlation and Multicollinearity Boundaries:}
    \begin{itemize}
        \item Operating profitability (ROA) correlated almost perfectly with EPS ($r = 0.975, p < 0.01$), confirming that simultaneous inclusion would generate severe multicollinearity (variance inflation factor (VIF) $\approx 20$). Hence, ROA was preserved strictly as an isolated substitute.
        \item EPS and BVPS have a moderate bivariate correlation ($r = 0.372, p = 0.051$). Model~1 substantive-regressor VIFs are approximately 1.16. Model~2 VIFs are 10.33 (EPS), 2.28 (BVPS), 15.15 (EPS $\times$ Post), and 24.13 (BVPS $\times$ Post), warranting caution in interpreting individual coefficients. On the full fitted design the VIFs are 5.70 and 12.32 (Model~1) and 14.82, 12.46, 27.34 and 259.32 (Model~2).
    \end{itemize}
    \item \textbf{Baseline Value Relevance (Model 1):}
    \begin{itemize}
        \item In the baseline two-way fixed effects model, the estimated earnings coefficient was $\hat{\beta}_E = 2.095$ (HC3 SE = 4.096, $t = 0.51, p = 0.616$ on $t(16)$) and the book value coefficient was $\hat{\beta}_B = 0.219$ (HC3 SE = 1.858, $t = 0.12, p = 0.908$ on $t(16)$).
        \item Both coefficients were positive, consistent with valuation theory, but neither was statistically significant at the 5\% level. The imprecision comes from the bank fixed effects, not the year effects: with year effects only, the earnings slope is 2.185 (SE 1.260, $p = 0.098$), and adding bank effects raises its standard error to 4.096, so the bank effects absorb the cross-bank variation that identifies the slope. The model yielded $R^2 = 0.948$ (Adjusted $R^2 = 0.912$). The sharp increase from $R^2 = 0.136$ in pooled ordinary least squares (OLS) to $0.919$ with year effects alone ($0.948$ with bank effects added) illustrates how much price variation common fiscal-year effects absorb.
    \end{itemize}
    \item \textbf{Regime Shift and Implied-Slope Compression (Model 2):}
    \begin{itemize}
        \item \textit{Pre-Consolidation Valuation:} The earnings slope estimate was $\hat{\beta}_E = 5.577$ (HC3 SE = 3.560, $p = 0.140$), and the book value slope estimate was $\hat{\beta}_B = 0.100$ (HC3 SE = 1.306, $p = 0.940$). The earnings point estimate is larger, but both estimates are imprecise, and their relative magnitude does not establish which measure was more value relevant.
        \item \textit{Regime Interaction Shifts:} The post-consolidation interaction terms were negative: $\hat{\delta}_E = -3.124$ (HC3 SE = 4.003, $t = -0.78, p = 0.448$) and $\hat{\delta}_B = -1.758$ (HC3 SE = 1.231, $t = -1.43, p = 0.175$).
        \item \textit{Implied Post-Consolidation Slopes:} The implied post-consolidation point estimates are $\hat{\beta}_E + \hat{\delta}_E = 2.453$ ($p = 0.298$) and $\hat{\beta}_B + \hat{\delta}_B = -1.658$ ($p = 0.161$). These estimates are imprecise and do not establish a decline in either slope or a negative book-value association.
        \item \textit{Principal Joint Wald Test:} The joint test of parameter stability ($H_0: \delta_E = \delta_B = 0$) yielded $F(2, 14) = 2.954$ ($p = 0.085$) under the pre-specified HC3 covariance. Under the designated 5\% threshold, the baseline model does not reject parameter stability under HC3, though it achieves marginal significance at the 10\% level. The same test rejects at 5\% under the classical, HC0, HC1 and HC2 covariances for all eight specifications R2--R9 (Chapter IV, Table~\ref{tab:cov_sensitivity}), so the baseline non-rejection is specific to the HC3 estimator.
    \end{itemize}
    \item \textbf{Robustness Suite Insights and the Entity-Timing Sensitivity:}
    \begin{itemize}
        \item The robustness suite (R1--R10) did not reject parameter stability at the 5\% level across ROA substitution (R3), cash dividend inclusion (R4, residual $\text{df} = 13$), transition-year removal (R5), bonus share dilution (R7), restated basic earnings (R8), Kanchan-mismatched bank-year removal (R9, residual $\text{df} = 13$), and supplementary partial-year Deva removal (R10, residual $\text{df} = 12$), except under Specification R6, whose unadjusted joint Wald test rejected at 5\%. This is the only raw rejection among nine joint tests R2--R10 under HC3, and the Holm-adjusted p-value is 0.190 ($p = 0.169$ across R2--R9).
        \item \textit{Entity-Timing Sensitivity (Specification R6):} When FY 2078/79 was excluded to address the entity-timing asymmetry (standalone pre-merger fundamentals paired with April 2023 trading prices for the four merged banks), the unadjusted joint test was below 5\%:
        \begin{equation}
        F(2, 10) = 5.82, \quad p = 0.021 \quad (N = 23)
        \end{equation}
        with $\hat{\delta}_B = -2.611$ (HC3 SE = 1.517, $t = -1.72, p = 0.116$ on $t(10)$). Neither interaction is individually significant. R5, which has the same sample size and residual degrees of freedom, does not reject (F(2,10)=2.41, p=0.139). R6 is an unadjusted sensitivity result, and its Holm-adjusted p-value is 0.190 ($p = 0.169$ across R2--R9). In the R6 sample, dropping any single bank leaves the joint test far from rejection ($F = 0.615$--$1.309$, $p = 0.329$--$0.572$), so the R6 rejection does not survive a leave-one-bank perturbation.
    \end{itemize}
    \item \textbf{Specification Comparisons and Omitted Variable Bias:}
    Pooled OLS produced an EPS slope estimate of $4.882$ ($R^2=0.136$), compared with $2.095$ in the two-way fixed-effects model ($R^2=0.948$). The pooled estimate is imprecise ($p=0.128$), and the raw $R^2$ values include the variation accounted for by the fixed effects. They do not isolate the contribution of the accounting regressors.
\end{itemize}

No prospective power analysis was specified. A post hoc design-sensitivity calculation using two-sided $\alpha=0.05$, 80\% power, fitted HC3 standard errors, and the residual-$df$ $t$ reference gives approximate detectable effects of 10.73 NPR for the Model~2 pre-period EPS slope, 12.06 NPR for the EPS interaction, and 3.71 NPR for the BVPS interaction. These thresholds show that the non-rejections do not rule out effects smaller than those magnitudes. They are not observed-power estimates.

\section{Conclusions}

Based on the empirical findings, econometric diagnostics, and theoretical evaluations conducted in this dissertation, the following major conclusions, direct answers to research questions, and formal hypothesis decisions are drawn:

\subsection{Direct Answers to Research Questions}
In direct response to the three research questions formulated in Chapter~I:
\begin{itemize}
    \item \textbf{Answer to Research Question 1 (Baseline Value Relevance):} The baseline two-way fixed-effects model does not find statistically significant EPS or BVPS associations with the price measure ($p = 0.616$ for EPS, and $p = 0.908$ for BVPS). Including year fixed effects accounts for common macro-year price variation and raises model $R^2$. This comparison does not establish that omitted year shocks account for the explanatory power of pooled studies.
    \item \textbf{Answer to Research Question 2 (Consolidation Regime Shift):} The estimated interaction coefficients are negative ($\hat{\delta}_E=-3.124$, $\hat{\delta}_B=-1.758$), but neither is statistically significant. The baseline joint test does not reject parameter stability at 5\% ($F=2.954$, $p=0.085$). The point estimates alone do not establish downward compression.
    \item \textbf{Answer to Research Question 3 (Sensitivity and Boundary Conditions):} The unadjusted R6 joint test is below 5\% ($F(2,10)=5.82$, $p=0.021$), but this is the only raw rejection among nine tests R2--R10 under HC3. Its Holm-adjusted $p$-value is 0.190 ($p = 0.169$ across R2--R9). The remaining specifications do not reject at 5\% under HC3. The evidence is sensitivity-dependent.
\end{itemize}

\subsection{Regime Differences in Accounting Valuation Slopes}
The pre-consolidation point estimates were $\hat{\beta}_E=5.58$ and $\hat{\beta}_B=0.10$, whereas the implied post-consolidation point estimates were 2.45 and $-1.66$. These estimates are imprecise and do not establish earnings dominance, a decline in either slope, or a negative book-value association. The baseline joint test does not reject at 5\% ($F=2.954$, $p=0.085$). The R6 result is an unadjusted sensitivity finding and is not significant after Holm adjustment across the nine joint tests (adjusted $p=0.190$, with $p = 0.169$ across R2--R9). Results are reported as associations across calendar periods, not as an established causal merger impact.

\subsection{Behavior of Book Value and Earnings Slopes Across Regimes}
Under classical valuation theory \citep{ohlson1995,collins1997}, disrupted earnings may increase the valuation role of book value. The negative post-consolidation BVPS point estimate ($\hat{\beta}_B + \hat{\delta}_B=-1.658$) is statistically insignificant ($p=0.161$), so these data do not establish whether this theoretical substitution occurred or how market participants treated expanded book equity.

\subsection{Econometric Discipline: Specification Sensitivity and Macro-Year Variation}
This dissertation highlights the influence of econometric specification in evaluating accounting value relevance. Pooled and two-way fixed-effects estimates differ, but raw $R^2$ includes variation absorbed by fixed effects and should not be interpreted as the contribution of accounting fundamentals alone. Slope estimates are conditional on both bank and fiscal-year fixed effects. They are not obtained by simple double demeaning. P/E is excluded because it embeds the dependent price in its numerator. The estimated accounting associations remain uncertain in this small sample.

\subsection{Archival Timing and Data Integrity in Frontier Market Research}
The sensitivity analysis shows dependence on sample and entity-boundary definitions. Excluding FY 2078/79 in R6 yielded an unadjusted joint test of $F(2,10)=5.82$ ($p=0.021$), but across the nine joint tests R2--R10 its Holm-adjusted $p$-value is 0.190 ($p = 0.169$ across R2--R9). The observed F is 41.8\% above the 5\% critical value. The five-day cutoff is an operational rule. Including the two delayed observations changes the baseline joint test to $F(2,16)=4.762$ ($p=0.024$). Carrying forward the 12 May 2021 panel-series prices for HBL/NIB ($487/453$) gives $F(2,16)=3.870$ ($p=0.043$), while using the alternative 13 May reported pair ($484/460$) gives $F(2,16)=3.821$ ($p=0.044$). Both pairs are roughly a year before the April 2022 target and are stale-price scenarios, not observations satisfying the freshness rule. In the $N=30$ model with the observed delayed prices, the pre-period BVPS estimate is 0.709 compared with 0.100 in the baseline. These sensitivities show that the sample rule is consequential. A wild-cluster bootstrap of the two joint Wald tests, enumerating all 32 Rademacher sign patterns at five banks, gives $p = 0.0625$ for both the baseline and R6, flipping only the R6 decision on a single pattern. With 32 patterns the smallest attainable $p$-value is $2/32 = 0.0625$, so the flip is mechanically forced, making the check inconclusive at this cluster count and leaving the HC3 inference unchanged.

\subsection{Formal Decisions on Research Hypotheses}
To formally synthesize the empirical findings against the theoretical expectations formulated in Chapter~I, Table~\ref{tab:hyp_decisions} presents the decisions on the four operational research hypotheses and the joint structural stability hypothesis.

\begin{table}[!htbp]
\centering
\small
\caption{Summary of Hypothesis Testing Decisions}
\label{tab:hyp_decisions}
\begin{tabular}{>{\raggedright\arraybackslash}p{0.21\textwidth}>{\raggedright\arraybackslash}p{0.17\textwidth}>{\raggedright\arraybackslash}p{0.14\textwidth}>{\raggedright\arraybackslash}p{0.19\textwidth}>{\raggedright\arraybackslash}p{0.14\textwidth}}
\toprule
\textbf{Research Hypothesis} & \textbf{Model / Parameter} & \textbf{Null Hypothesis} & \textbf{Estimate (SE) / Stat.} & \textbf{Empirical Decision} \\
\midrule
$H_1$: EPS Relevance & Model 1 Baseline $\beta_E$ & $H_0: \beta_E = 0$ & 2.095 (4.096), $p = 0.616$ & Fail to Reject $H_0$ \\
\addlinespace
& Model 2 Pre $\beta_E$ & $H_0: \beta_E = 0$ & 5.577 (3.560), $p = 0.140$ & Fail to Reject $H_0$ \\
\addlinespace
$H_2$: BVPS Relevance & Model 1 Baseline $\beta_B$ & $H_0: \beta_B = 0$ & 0.219 (1.858), $p = 0.908$ & Fail to Reject $H_0$ \\
\addlinespace
& Model 2 Pre $\beta_B$ & $H_0: \beta_B = 0$ & 0.100 (1.306), $p = 0.940$ & Fail to Reject $H_0$ \\
\addlinespace
$H_3$: EPS Regime Shift & Model 2 Interaction $\delta_E$ & $H_0: \delta_E = 0$ & $-3.124$ (4.003), $p = 0.448$ & Fail to Reject $H_0$ \\
\addlinespace
$H_4$: BVPS Regime Shift & Model 2 Interaction $\delta_B$ & $H_0: \delta_B = 0$ & $-1.758$ (1.231), $p = 0.175$ & Fail to Reject $H_0$ \\
\addlinespace
Joint Regime Shift & Model 2 Joint Wald & $H_0: \delta_E=\delta_B=0$ & $F(2, 14) = 2.954, p = 0.085$ & Fail to Reject $H_0$ \\
\addlinespace
(Composite Test) & R6 Joint Wald (Sens.) & $H_0: \delta_E=\delta_B=0$ & $F(2, 10) = 5.818, p = 0.021$ & Sensitivity check only, no decision (Holm-adjusted $p = 0.190$) \\
\bottomrule
\end{tabular}

\vspace{4pt}
\begin{minipage}{\linewidth}
\footnotesize
\textit{Note.} Decisions for the $t$-tests evaluated at $\alpha = 0.05$ two-tailed, whereas the joint Wald $F$-tests are evaluated against the upper tail. SE denotes standard error. Standard errors and Wald tests estimated under MacKinnon-White HC3 robust covariance matrix on Student's $t$ residual degrees of freedom ($\text{df} = 16$ for Model~1, $\text{df} = 14$ for Model~2, $\text{df} = 10$ for R6). \\
\textit{Abbreviations:} EPS (Earnings Per Share), BVPS (Book Value Per Share), $\beta_E$ and $\beta_B$ (valuation slopes), $\delta_E$ and $\delta_B$ (consolidation interaction shifts), SE (Standard Error), df (degrees of freedom), HC3 (MacKinnon-White heteroskedasticity-consistent covariance estimator).
\end{minipage}
\end{table}

As summarized in Table~\ref{tab:hyp_decisions}, the empirical evaluations lead to the following formal decisions:
\begin{itemize}
    \item \textbf{Hypothesis 1 ($H_1$: Value Relevance of EPS):} Evaluated via the null hypothesis $H_0: \beta_E = 0$. In the baseline two-way fixed effects Model~1, the earnings slope was $\hat{\beta}_E = 2.095$ ($t = 0.51, p = 0.616$ on $t(16)$). In Model~2, the pre-consolidation earnings slope was $\hat{\beta}_E = 5.577$ ($t = 1.57, p = 0.140$ on $t(14)$). Because both $p$-values exceed the 5\% significance threshold, the study fails to reject the null hypothesis $H_0$. Accounting earnings do not exhibit a statistically significant association with post-reporting-deadline market prices in the two-way fixed-effects specification, where bank effects absorb the cross-bank variation that identifies the slope.
    \item \textbf{Hypothesis 2 ($H_2$: Value Relevance of BVPS):} Evaluated via the null hypothesis $H_0: \beta_B = 0$. In Model~1, the book value slope was $\hat{\beta}_B = 0.219$ ($t = 0.12, p = 0.908$ on $t(16)$). In Model~2, the pre-consolidation book value slope was $\hat{\beta}_B = 0.100$ ($t = 0.08, p = 0.940$ on $t(14)$). Failing to attain statistical significance at the 5\% level, the evidence fails to reject $H_0$. Book value per share shows no statistically significant incremental pricing effect in the baseline specification.
    \item \textbf{Hypothesis 3 ($H_3$: Structural Regime Shift in EPS Slope):} Evaluated via the null hypothesis $H_0: \delta_E = 0$ in Model~2. The estimated interaction parameter was $\hat{\delta}_E = -3.124$ ($t = -0.78, p = 0.448$ on $t(14)$). Because $p > 0.05$, the study fails to reject $H_0$. The individual compression in the earnings valuation slope across the 2023 consolidation wave is not statistically significant.
    \item \textbf{Hypothesis 4 ($H_4$: Structural Regime Shift in BVPS Slope):} Evaluated via the null hypothesis $H_0: \delta_B = 0$ in Model~2. The estimated interaction parameter was $\hat{\delta}_B = -1.758$ ($t = -1.43, p = 0.175$ on $t(14)$). Because $p > 0.05$, the study fails to reject $H_0$. The individual shift in the book value valuation slope is not statistically significant at the 5\% level.
    \item \textbf{Joint Test of Structural Regime Shift:} Evaluated via the joint null hypothesis $H_0: \delta_E = \delta_B = 0$. In baseline Model~2, the Wald test yielded $F(2,14)=2.954$ ($p=0.085$), failing to reject at 5\%. Sensitivity check R6 yielded an unadjusted $F(2,10)=5.82$ ($p=0.021$), which is the only raw rejection among nine tests R2--R10 under HC3, with a Holm-adjusted $p$-value of 0.190 ($p = 0.169$ across R2--R9). The evidence for a regime shift is therefore sensitivity-dependent and not conclusive.
\end{itemize}

Two limits stated in Chapter~I bear on how far these conclusions travel. First, this purposive five-series sample includes four selected combinations from the early-2023 consolidation wave and one unmerged benchmark. The fifth combination, Prabhu Bank/Century, was not selected, and no separate data-based exclusion criterion is claimed. The findings therefore describe the sampled institutions and do not establish results for the omitted Prabhu/Century combination or the full commercial banking population. They cannot be generalized to development banks, finance companies, or non-financial listed sectors. Second, actual publication dates were verified in the reviewed records for only 2 of the 30 bank-years. For the remaining observations, the statutory deadline does not establish when the information became public, so the post-reporting-deadline pricing protocol rests on the deadline as an operational rule rather than on verified publication in each case.

\section{Implications of the Study}

The findings and methodological insights of this dissertation offer targeted implications partitioned across three institutional domains:

\subsection{Implications for Regulatory Authorities (NRB, SEBON, and NEPSE)}
\begin{itemize}
    \item \textbf{Scrutiny of Post-Merger Asset Quality and Provisioning Transparency (NRB):}
    The imprecise point estimates suggest a downward compression in the earnings slope and no positive book value pricing anchor post-consolidation. Whether this coincides with market caution regarding the asset quality of merged commercial banks is a hypothesis: no asset-quality variable is estimated in this study. Given that uncertainty, NRB's supervisory departments could prioritize thorough on-site asset quality reviews (AQRs) across consolidated entities. Standardized, transparent impairment testing of loan books assimilated during amalgamations is essential to restore market confidence in reported financial accounts.
    \item \textbf{Evaluation of Mandatory Consolidation Policies (NRB):}
    While consolidation mandates reduced the licensed commercial bank population from 32 in FY 2011/12 to 20 by mid-October 2023, the imprecise estimates do not establish that horizontal amalgamations were accompanied by substantial earnings transitoriness (no transitoriness parameter was estimated) or by an attenuation of the valuation slopes, since neither interaction is statistically significant. Policy frameworks should not assume that structural combinations automatically translate into improved financial stability or higher equity valuations in the short to medium term, because financial stability was not measured in this study.
    \item \textbf{Transparency in Capital Distribution Directives (NRB):}
    The sharp reduction in cash dividend distributions post-consolidation removed an observable cash flow signal for bank equities, though DPS was never shown to anchor valuation (its correlation with price is $r = 0.082$, $p = 0.679$). NRB should provide clear, forward-looking guidelines regarding dividend distribution criteria and counter-cyclical capital buffer expectations, helping market participants formulate realistic cash flow projections.
    \item \textbf{Reforming Secondary Market Trading Suspension Practices (SEBON and NEPSE):}
    A critical institutional friction identified in this study was the extended merger-related trading suspensions, which froze secondary market trading for up to 54 days past statutory reporting deadlines. SEBON and NEPSE should enforce trading continuity during corporate combinations, restricting trading halts to at most a few sessions around share swap execution. Continuous price discovery is necessary to maintain liquidity and avoid information distortion.
    \item \textbf{Standardizing Corporate Disclosure Timelines and Digital Access (SEBON and NEPSE):}
    Under Section 59(2) of BAFIA, commercial banks are legally mandated to publish audited financial statements within nine months into the subsequent fiscal year \citep{bafia2017}. To mitigate the risk of fragmented corporate disclosures and ensure uniform public access, SEBON and NEPSE should operationalize a centralized electronic disclosure portal where audited annual reports, quarterly Nepal Financial Reporting Standards (NFRS) financials, and merger swap audit reports are uploaded instantaneously with certified timestamps.
    \item \textbf{Discouraging Endogenous Ratios in Equity Research (SEBON and NEPSE):}
    Given the historical use of endogenous price-earnings multiples in local equity research and academic literature, regulatory education frameworks should promote sound financial modelling practices that avoid regressing share price on price-derived ratios.
    \item \textbf{Standardizing NFRS Restatement Disclosures (Accounting Standard Setters):}
    The Accounting Standards Board (ASB) of Nepal and the Institute of Chartered Accountants of Nepal (ICAN) could consider enforcing uniform reporting of bonus-adjusted basic EPS and diluted EPS across consecutive annual reports. This study does not establish that current reporting is non-uniform, so the recommendation rests on transparency grounds rather than on a documented deficiency.
\end{itemize}

\subsection{Implications for Commercial Bank Executives and Boards of Directors}
\begin{itemize}
    \item \textbf{Prioritizing Post-Merger Operational and Credit Integration:}
    The observed earnings dispersion post-consolidation is a descriptive change whose cause is not identified by this study. To the extent that dispersion reflects integration friction, executive leadership may benefit from focusing on internal operational harmonization, including the integration of management information systems, core banking software, and credit underwriting standards. The results do not establish that these specific actions are necessary.
    \item \textbf{Transparent Disclosure of Restructuring Adjustments:}
    Given the uncertainty around the post-consolidation book-value association, commercial bank management should enhance voluntary footnote disclosures regarding purchase price allocations, recognized goodwill, restructuring provisions, and reserve transfers. Granular disclosures regarding asset quality can help alleviate investor skepticism.
    \item \textbf{Communication of Sustainable Capital Policies:}
    Bank executives should pursue sustainable capital management strategies that balance regulatory capital adequacy maintenance with transparent, predictable dividend distribution policies.
\end{itemize}

\subsection{Implications for Equity Investors, Financial Analysts, and Market Participants}
\begin{itemize}
    \item \textbf{Scrutinizing Book Value Composition:}
    Equity investors and analysts in Nepal should exercise caution when using unadjusted book value per share as a safety metric for consolidated banks. In the post-merger banking era, reported book equity may encompass capital issued against stressed assets or restructuring costs, though no balance-sheet composition variable was estimated to establish this. Fundamental analysis should examine the underlying quality of net assets, adjusting for unamortized goodwill and loan provisioning levels.
    \item \textbf{Recognizing Earnings Transitoriness in Merged Institutions:}
    Market participants should account for the fact that reported accounting earnings in the immediate years following major amalgamations can exhibit heightened volatility. Whether that volatility reflects transitoriness due to credit-loss charges and integration expenses was not measured. Valuation models should use normalized operational metrics rather than headline basic EPS.
    \item \textbf{Incorporating Macro-Temporal Factors:}
    Financial analysts must recognize that common fiscal-year shocks exert a dominant influence on bank share prices (year fixed effects raise $R^2$ from 0.136 to 0.919). The specific channels (systemic liquidity, policy rates, and broad market sentiment) are not identified by this study, since no macroeconomic variable is in any model. Valuation frameworks should incorporate macro-temporal factors rather than relying solely on static accounting slopes.
\end{itemize}

\section{Recommendations and Directions for Future Research}

\subsection{Policy and Managerial Recommendations}
Based on the synthesized conclusions and institutional implications, the following specific recommendations are offered:
\begin{itemize}
    \item \textbf{Mandatory Post-Merger Asset Quality Disclosures:} Given that the recovery performance, provisioning coverage, and write-offs of acquired legacy loan portfolios are unmeasured in this study and therefore an acknowledged evidence gap, Nepal Rastra Bank could consider mandating that commercial banks completing horizontal amalgamations publish a semi-annual ``Asset Quality and Integration Disclosure Report'' for three years post-merger, detailing those items.
    \item \textbf{Limiting Trading Suspensions During Corporate Restructuring:} The Securities Board of Nepal and NEPSE should establish regulatory ceilings that restrict secondary market trading halts during corporate combinations to at most a few consecutive trading sessions around formal share swap execution.
    \item \textbf{Operationalizing a Centralized Digital Filings Repository:} SEBON should establish an open-access, machine-readable digital filings portal requiring all listed financial institutions to upload audited financial statements, board resolutions, and statutory disclosures concurrently with their public release.
    \item \textbf{Publishing Core Operational Earnings Metrics:} Commercial bank management and equity research institutions should report normalized operating earnings per share alongside statutory NFRS basic EPS, distinguishing ongoing operational performance from non-recurring merger adjustments and regulatory reserve transfers. No non-recurring merger adjustment or regulatory reserve transfer is quantified in this study, so the recommendation is a transparency proposal, not a response to a measured deficiency.
\end{itemize}

\subsection{Avenues for Future Academic Inquiry}
While this dissertation has established an econometric baseline for evaluating accounting value relevance across the 2023 consolidation wave, several fruitful avenues remain open for future scholarly research:
\begin{itemize}
    \item \textbf{Methodological Rigor and Econometric Specification in Frontier Markets:}
    This dissertation demonstrates the importance of sound panel econometric specifications in Nepalese accounting research. Graduate researchers and academic institutions should eliminate mathematically circular regressors (such as P/E predicting MPS), incorporate two-way fixed effects to control for unobserved bank and temporal heterogeneity, and employ finite-sample robust inference methods such as MacKinnon-White HC3 on Student's $t$ residual degrees of freedom.
    \item \textbf{Extension to Long-Horizon Post-Consolidation Panels:}
    This study examined the post-consolidation regime across three initial fiscal years (FY 2079/80 through FY 2081/82). As additional financial years elapse, future research should re-estimate the interactive valuation models over longer post-merger horizons (5 to 10 years) to evaluate whether valuation slopes change once integration is fully completed.
    \item \textbf{Expansion to the Full Commercial Banking Population:}
    Future studies should expand the cross-sectional sample to encompass all 19 NEPSE-listed commercial banks, employing difference-in-differences (DiD) panel frameworks that compare aggressively merged banking series against a cohort of organically grown or minimally consolidated peer institutions, if such a cohort can be identified among the 19 listed banks.
    \item \textbf{Investigation of Non-Financial Disclosure and Environmental, Social, and Governance (ESG) Metrics:}
    Subsequent research should examine the value relevance of non-financial indicators, corporate governance metrics, digital banking investments, and Environmental, Social, and Governance (ESG) disclosures among Nepalese financial institutions, assessing whether non-accounting disclosures provide incremental pricing power.
    \item \textbf{Intraday Event-Study Methodologies:}
    Building upon the panel findings of this study, future inquiries could employ high-frequency, intraday event-study methodologies to capture the immediate market reaction surrounding statutory board dividend announcements, AGM approvals, and quarterly earnings releases, quantifying the speed and informational efficiency of NEPSE's price discovery mechanism.
    \item \textbf{Cross-Sectoral Valuation Comparative Inquiries:}
    Finally, comparative studies should investigate accounting value relevance across divergent sectors of the Nepalese economy, contrasting the commercial banking industry with insurance companies, hydropower developers, microfinance institutions, and manufacturing corporations to evaluate how divergent regulatory environments shape the pricing of earnings and book value.
\end{itemize}

